% Source: https://artofproblemsolving.com/community/c3h2460329_please_help_with_this_problem
% Stable identity: aops-topic-2460329-q01
Number the squares $0,1,2,\dots$, square $k$ having side $\dfrac1{2^k}$ and area $\dfrac1{4^k}$. Square $k$ ($k\ge1$) is bisected by the top side of square $k-1$: its lower half lies inside square $k-1$, and its upper half lies above the top side of square $k-1$. The top sides of the squares get higher and higher, so the top side of square $k-1$ is the highest point of squares $0,1,\dots,k-1$; hence the upper half of square $k$ does not overlap any earlier square. So each new square adds exactly half its area:
\[1+\dfrac12\left(\dfrac14+\dfrac1{16}+\dfrac1{64}+\cdots\right)=1+\dfrac12\cdot\dfrac{\dfrac14}{1-\dfrac14}=1+\dfrac12\cdot\dfrac13=\boxed{\dfrac76}.\]

\begin{asy}
usepackage("amsmath");
size(3.5cm);
defaultpen(fontsize(10pt));
real y=0, s=1; real c=0.5;
fill(box((0,0),(1,1)),rgb(0.85,0.85,0.85));
for(int k=1;k<=6;++k){
  real t=s/2; real top=y+s;
  fill(box((c-t/2,top),(c+t/2,top+t/2)),rgb(0.8,0.88,1));
  draw(box((c-t/2,top-t/2),(c+t/2,top+t/2)));
  y=top-t/2; s=t;
}
draw(box((0,0),(1,1)),linewidth(0.8));
label("$1$",(0.5,0),S);
draw((0.24,1.1)--(-0.05,1.1),gray); label("$\dfrac12$",(-0.05,1.1),W,fontsize(8pt));
draw((0.63,1.3)--(1.0,1.33),gray); label("$\dfrac14$",(1.0,1.33),E,fontsize(8pt));
\end{asy}

(A common mistake is to treat the added areas $\dfrac18,\dfrac1{32},\dfrac1{128},\dots$ as having ratio $\dfrac18$, which gives $1+\dfrac{\dfrac18}{1-\dfrac18}=\dfrac87$; the ratio is $\dfrac14$, because each square has a quarter of the area of the previous one.)
